\subsubsection{共振窗内 Hankel-null 的整系数主理想化与有限确定性}\label{subsubsec:pom-resonance-hankel-null-integral-principalization}
在共振窗 \(q=9,\dots,17\) 中，最小递推特征多项式 \(P_q(\lambda)\) 与 Hankel-null 证书之间存在一个整系数层面的主理想化对应：审计窗口内出现的任何整系数消去关系都必然来自 \(P_q\) 的短倍数，因此不存在“截断伪消去”。
该结构不仅给出核空间的完全刻画，也使得 \(P_q\) 与 Perron 根 \(r_q\) 可由有限 Hankel 核证书唯一恢复。

\begin{definition}[共振窗 Hankel-null 与短倍乘子模]\label{def:pom-resonance-hankel-null-multiples}
对 \(q\in\{9,\dots,17\}\)，令 \(h=\lfloor q/2\rfloor\)、\(n=2h+1\)，并令
\[
a_t:=S_q(t+2)\qquad(t\ge 0).
\]
取审计窗口长度 \(L\ge 1\)，定义矩阵 \(H=H(q)\in\ZZ^{n\times L}\) 为
\[
H_{i+1,j+1}:=a_{i+j}\qquad(0\le i\le n-1,\ 0\le j\le L-1),
\]
并定义
\[
\mathrm{NullModes}(q):=\ker_{\ZZ}(H^{\mathsf T})\subset\ZZ^{n}.
\]
令 \(P_q(\lambda)\in\ZZ[\lambda]\) 表示序列 \(t\mapsto a_t\) 的首一最小湮灭多项式（等价于 \(m\mapsto S_q(m)\) 的最小整系数递推特征多项式；定理 \ref{thm:pom-moment-resonance}），记 \(d_q:=\deg P_q\) 与缺口 \(\Delta(q):=n-d_q\)。
定义截断倍乘子模
\[
\mathcal{M}(P_q):=
\Bigl\{\operatorname{coeff}_{<n}\!\bigl(Q(\lambda)P_q(\lambda)\bigr)\in\ZZ^n:\ 
Q\in\ZZ[\lambda],\ \deg Q<\Delta(q)\Bigr\}\subset\ZZ^n.
\]
\end{definition}

\begin{theorem}[共振窗内的整系数主理想化]\label{thm:pom-resonance-hankel-null-integral-principalization}
\leanverified{paper\_pom\_resonance\_hankel\_null\_integral\_principalization}
在定义 \ref{def:pom-resonance-hankel-null-multiples} 的设定下，若审计窗口满足 \(L\ge n\)，则在 \(\ZZ\)-模意义下成立严格恒等式
\[
\boxed{\;\mathrm{NullModes}(q)=\mathcal{M}(P_q).\;}
\]
\end{theorem}
\begin{proof}
对序列 \((a_t)_{t\ge 0}\) 应用定理 \ref{thm:pom-hankel-syndrome-module-kernel-equals-multiples}，取 \(m_0=0\) 并令其方块 Hankel 矩阵阶数为 \(n\)。
由于 \(L\ge n\)，条件 \(H^{\mathsf T}\alpha=0\) 等价于对所有 \(0\le k\le n-1\) 有
\(
\sum_{i=0}^{n-1}\alpha_i a_{k+i}=0
\)，从而与方块 Hankel 的转置核一致。
又由 \(P_q\) 的首一最小性假设，上述定理给出
\(
\ker_{\ZZ}(H^{\mathsf T})=\mathcal{M}(P_q)
\)，即得结论。证毕。
\end{proof}

\begin{corollary}[无截断伪消去与短倍数分解的唯一性]\label{cor:pom-resonance-no-spurious-elimination-unique-short-multiple}
\leanverified{paper\_pom\_resonance\_no\_spurious\_elimination\_unique\_short\_multiple}
在定理 \ref{thm:pom-resonance-hankel-null-integral-principalization} 的设定下，令
\(
B(\lambda)\in\ZZ[\lambda]
\)
满足 \(\deg B<n\)，且其系数向量 \(\operatorname{coeff}_{<n}(B)\in \mathrm{NullModes}(q)\)。
则存在唯一的 \(Q(\lambda)\in\ZZ[\lambda]\)（\(\deg Q<\Delta(q)\)）使得
\[
\boxed{\;B(\lambda)=Q(\lambda)P_q(\lambda).\;}
\]
因此任何在审计 Hankel 窗内成立的整系数消去关系必为最小递推特征多项式的短倍数，并自动延拓为全局递推关系。
\end{corollary}
\begin{proof}
由定理 \ref{thm:pom-resonance-hankel-null-integral-principalization} 有 \(\operatorname{coeff}_{<n}(B)\in\mathcal{M}(P_q)\)，故存在 \(Q\) 使 \(B=QP_q\) 且 \(\deg Q<\Delta(q)\)，得到存在性。
由于 \(P_q\) 首一且 \(\deg(QP_q)\le (\Delta(q)-1)+d_q=n-1<n\)，截断映射在次数 \(<n\) 的多项式上为同构；从而 \(Q\) 唯一。
最后，“延拓”为 \(P_q(E)a\equiv 0\) 的直接推论。证毕。
\end{proof}

\begin{corollary}[有限 Hankel 核证书的最大公因子恢复]\label{cor:pom-resonance-gcd-recovers-minpoly}
在定理 \ref{thm:pom-resonance-hankel-null-integral-principalization} 的设定下，取 \(\mathrm{NullModes}(q)\) 的任一 \(\ZZ\)-基
\(
\{\alpha^{(1)},\dots,\alpha^{(\Delta(q))}\}
\)
并定义相应湮灭多项式
\[
A^{(j)}(\lambda):=\sum_{k=0}^{n-1}\alpha^{(j)}_k\lambda^k\in\ZZ[\lambda].
\]
则在 \(\ZZ[\lambda]\) 中有
\[
\boxed{\;\gcd\bigl(A^{(1)}(\lambda),\dots,A^{(\Delta(q))}(\lambda)\bigr)=\pm P_q(\lambda).\;}
\]
从而 \(P_q\)（取首一规范化后）可由有限 Hankel 核证书唯一恢复；其 Perron 根 \(r_q\) 亦随之被完全确定。
\end{corollary}
\begin{proof}
由定理 \ref{thm:pom-resonance-hankel-null-integral-principalization}，每个 \(A^{(j)}\) 都可写为 \(A^{(j)}=Q^{(j)}P_q\)（\(\deg Q^{(j)}<\Delta(q)\)）。
另一方面，由 \(Q\equiv 1\) 的可容纳性可知 \(P_q\in\mathcal{M}(P_q)=\mathrm{NullModes}(q)\)，故存在整数 \(c_j\) 使
\[
P_q(\lambda)=\sum_{j=1}^{\Delta(q)} c_jA^{(j)}(\lambda)=\Bigl(\sum_{j=1}^{\Delta(q)}c_jQ^{(j)}(\lambda)\Bigr)P_q(\lambda).
\]
在整环 \(\ZZ[\lambda]\) 中推出 \(\sum_j c_jQ^{(j)}(\lambda)=1\)，因此 \(\{Q^{(j)}\}\) 生成单位理想，任意公因子必为 \(\pm1\)。
于是 \(\gcd_j A^{(j)}=\pm P_q\)。证毕。
\end{proof}

\paragraph{短整数证书、有限窗刚性与原始剩余的模素数扭转谱}
在上文的整系数主理想化之外，共振缺口还可从三个彼此兼容的有限量被完全压缩：短 NullMode 证书、有限窗核的主理想等式，以及原始剩余商的有限扭转谱。

\begin{theorem}[秩缺口强制短整数 NullMode 证书]\label{thm:pom-hankel-short-integer-certificate}
\leanverified{paper\_pom\_hankel\_short\_integer\_certificate}
设 \(A\in\ZZ_{\ge 0}^{R\times n}\) 满足
\(
\operatorname{rank}_{\QQ}(A)=r<n
\)，并令
\[
B:=\max_{i,j}A_{ij}\ge 1.
\]
则存在非零向量 \(x\in\ZZ^n\setminus\{0\}\) 使
\[
Ax=0,
\qquad
\|x\|_\infty\le (2nB)^{\frac{r}{n-r}}.
\]
因此一旦名义 \(n\)-维模态空间在整数矩阵口径下出现秩缺口 \(r<n\)，便必带有一条系数位长
\[
O\!\left(\frac{r}{n-r}\log(nB)\right)
\]
可控的整数零模证书。
\end{theorem}
\begin{proof}
从 \(A\) 中取一组 \(\QQ\)-线性无关的行，得到子矩阵
\(
A'\in \ZZ_{\ge 0}^{r\times n}
\)
且 \(\ker_{\QQ}(A')=\ker_{\QQ}(A)\)，故只需对 \(A'\) 证明结论。
以下仍记该 \(r\times n\) 矩阵为 \(A\)。

令
\[
Q:=(2nB)^{\frac{r}{n-r}},
\qquad
\mathcal S:=\{0,1,\dots,\lfloor Q\rfloor\}^{n}.
\]
则
\[
|\mathcal S|=(\lfloor Q\rfloor+1)^n>Q^n.
\]
对任意 \(u\in\mathcal S\)，由于 \(A\) 的条目非负且不超过 \(B\)，有
\[
0\le (Au)_i\le \sum_{j=1}^{n}A_{ij}\lfloor Q\rfloor\le nBQ
\qquad(1\le i\le r).
\]
故像集满足
\[
|A(\mathcal S)|\le (nBQ+1)^r\le (2nBQ)^r=Q^n.
\]
因此 \(|\mathcal S|>|A(\mathcal S)|\)，由抽屉原理存在 \(u\neq v\in\mathcal S\) 使 \(Au=Av\)。
令 \(x:=u-v\)，则 \(x\in\ZZ^n\setminus\{0\}\)、\(Ax=0\)，且
\[
\|x\|_\infty\le \lfloor Q\rfloor\le Q=(2nB)^{\frac{r}{n-r}}.
\]
证毕。
\end{proof}

\begin{corollary}[共振缺口的可控位长 NullMode 证书]\label{cor:pom-resonance-short-nullmode-certificate}
\leanverified{paper\_pom\_resonance\_short\_nullmode\_certificate}
沿用定义 \ref{def:pom-resonance-hankel-null-multiples} 的记号，并设 \(L\ge n\)。
若
\[
d_q=\deg P_q<n=2\Bigl\lfloor\frac q2\Bigr\rfloor+1,
\]
则存在非零向量 \(\alpha\in \mathrm{NullModes}(q)\subset\ZZ^n\) 使
\[
\|\alpha\|_\infty
\le
\bigl(2nB_q\bigr)^{\frac{d_q}{n-d_q}},
\qquad
B_q:=\max_{0\le i\le L-1,\ 0\le j\le n-1}a_{i+j},
\]
其中 \(a_t=S_q(t+2)\)。
特别地，共振缺口一旦出现，便必存在一条位长仅按
\(
O\!\bigl(\frac{d_q}{n-d_q}\log(nB_q)\bigr)
\)
增长的整数 Hankel-null 证书。
\end{corollary}
\begin{proof}
取 \(A:=H(q)^{\mathsf T}\in\ZZ_{\ge 0}^{L\times n}\)。
由于 \(L\ge n\) 且 \(\operatorname{rank}_{\QQ}(H(q))=d_q\)，有
\(
\operatorname{rank}_{\QQ}(A)=d_q<n
\)。
由定义 \(\mathrm{NullModes}(q)=\ker_{\ZZ}(H(q)^{\mathsf T})=\ker_{\ZZ}(A)\)。
将定理 \ref{thm:pom-hankel-short-integer-certificate} 应用于 \(A\) 即得结论。证毕。
\end{proof}

\begin{theorem}[有限窗 Hankel-null 的主理想刚性]\label{thm:pom-hankel-finite-window-principal-rigidity}
\leanverified{paper\_pom\_hankel\_finite\_window\_principal\_rigidity}
设 \(K\) 为域，\((a_m)_{m\ge 0}\subset K\) 为 Hankel 秩 \(d\) 的序列，并令 \(P(\lambda)\in K[\lambda]\) 为其首一最小递推特征多项式（故 \(\deg P=d\)）。
对整数 \(N\ge d\)、\(L\ge d+1\) 定义有限 Hankel 块
\[
\mathsf H_{N,L}:=(a_{i+j})_{0\le i\le N-1,\ 0\le j\le L-1}\in K^{N\times L}.
\]
若 \(\operatorname{rank}(\mathsf H_{N,L})=d\)，则在系数识别
\(
K^L\simeq K[\lambda]_{<L}
\)
下有严格等式
\[
\ker(\mathsf H_{N,L})
=
\mathcal M_L(P),
\]
其中
\[
\mathcal M_L(P):=
\Bigl\{
\operatorname{coeff}_{<L}\!\bigl(Q(\lambda)P(\lambda)\bigr):
\ Q\in K[\lambda],\ \deg Q\le L-d-1
\Bigr\}\subset K^L.
\]
换言之，在有限窗已经达到最小递推秩的前提下，Hankel-null 中不存在任何超出最小递推倍乘所强制的原始零模方向。
\end{theorem}
\begin{proof}
对任意 \(Q(\lambda)\in K[\lambda]\) 满足 \(\deg Q\le L-d-1\)，由于 \(P(E)a\equiv 0\)，故
\[
\bigl(QP\bigr)(E)a\equiv 0.
\]
因此 \(\operatorname{coeff}_{<L}(QP)\in\ker(\mathsf H_{N,L})\)，从而
\(
\mathcal M_L(P)\subseteq \ker(\mathsf H_{N,L})
\)。

另一方面，由 \(P\) 首一可知映射
\[
K[\lambda]_{\le L-d-1}\longrightarrow K^L,
\qquad
Q\longmapsto \operatorname{coeff}_{<L}(QP)
\]
为单射，故
\[
\dim_K\mathcal M_L(P)=L-d.
\]
又由 \(\operatorname{rank}(\mathsf H_{N,L})=d\)，有
\[
\dim_K\ker(\mathsf H_{N,L})=L-d.
\]
于是两者一包含且维数相同，故必相等。证毕。
\end{proof}

\begin{theorem}[整数原始剩余的扭转谱与坏素数集合]\label{thm:pom-hankel-primitive-torsion-spectrum}
在定理 \ref{thm:pom-hankel-finite-window-principal-rigidity} 的整数情形中，设 \((a_m)_{m\ge 0}\subset\ZZ\)、\(P(\lambda)\in\ZZ[\lambda]\) 为首一且本原的最小递推特征多项式，并令
\[
\operatorname{rank}_{\QQ}(\mathsf H_{N,L})=d.
\]
定义
\[
K_{\ZZ}:=\ker_{\ZZ}(\mathsf H_{N,L})\subset\ZZ^L,
\qquad
M_{\ZZ}(P):=
\Bigl\{
\operatorname{coeff}_{<L}\!\bigl(Q(\lambda)P(\lambda)\bigr):
\ Q\in\ZZ[\lambda],\ \deg Q\le L-d-1
\Bigr\}\subset\ZZ^L,
\]
以及原始剩余商
\[
T_{\mathrm{prim}}:=K_{\ZZ}/M_{\ZZ}(P).
\]
则：
\begin{enumerate}
  \item \(T_{\mathrm{prim}}\) 为有限阿贝尔群；
  \item 记
  \[
  \mathcal B_{\mathrm{rk}}(\mathsf H_{N,L})
  :=
  \left\{
  p\ \text{为素数}:\ \operatorname{rank}_{\FF_p}(\mathsf H_{N,L}\bmod p)<d
  \right\},
  \]
  并令 \(\overline{M}_p\subset(\FF_p)^L\) 表示 \(M_{\ZZ}(P)\) 的模 \(p\) 约化像。
  则对每个 \(p\notin \mathcal B_{\mathrm{rk}}(\mathsf H_{N,L})\) 有
  \[
  \dim_{\FF_p}
  \frac{\ker_{\FF_p}(\mathsf H_{N,L}\bmod p)}{\overline{M}_p}
  =
  \dim_{\FF_p}\!\bigl(T_{\mathrm{prim}}\otimes\FF_p\bigr).
  \]
  特别地，
  \[
  \ker_{\FF_p}(\mathsf H_{N,L}\bmod p)=\overline{M}_p
  \qquad\Longleftrightarrow\qquad
  p\nmid |T_{\mathrm{prim}}|.
  \]
\end{enumerate}
因此，模素数审计中出现“额外 Hankel-null 证书”的失真素数集合恰为
\[
\mathcal B_{\mathrm{rk}}(\mathsf H_{N,L})
\ \cup\
\{\,p:\ p\mid |T_{\mathrm{prim}}|\,\}.
\]
\end{theorem}
\begin{proof}
由定理 \ref{thm:pom-hankel-finite-window-principal-rigidity} 在 \(K=\QQ\) 上的结论，
\(
K_{\ZZ}\otimes_{\ZZ}\QQ=M_{\ZZ}(P)\otimes_{\ZZ}\QQ
\)
且两者秩均为 \(L-d\)。
故商 \(T_{\mathrm{prim}}=K_{\ZZ}/M_{\ZZ}(P)\) 只有扭转而无自由部分，从而为有限阿贝尔群，得第 (1) 条。

再记
\(
\overline{K}_p
\)
为 \(K_{\ZZ}\) 在 \((\FF_p)^L\) 中的模 \(p\) 约化像。
由于
\[
\ZZ^L/K_{\ZZ}\cong \operatorname{im}(\mathsf H_{N,L})\subset \ZZ^N
\]
为自由阿贝尔群，\(K_{\ZZ}\) 是 \(\ZZ^L\) 的饱和子模；故对每个素数 \(p\) 都有
\[
\dim_{\FF_p}\overline{K}_p=\operatorname{rank}_{\ZZ}(K_{\ZZ})=L-d.
\]
若 \(p\notin \mathcal B_{\mathrm{rk}}(\mathsf H_{N,L})\)，则
\[
\dim_{\FF_p}\ker_{\FF_p}(\mathsf H_{N,L}\bmod p)=L-d.
\]
而 \(\overline{K}_p\subseteq \ker_{\FF_p}(\mathsf H_{N,L}\bmod p)\)，维数又相同，故
\[
\overline{K}_p=\ker_{\FF_p}(\mathsf H_{N,L}\bmod p).
\]

现在对包含 \(M_{\ZZ}(P)\subset K_{\ZZ}\) 取 Smith 标准形。
存在 \(K_{\ZZ}\) 的一组基 \(\{e_1,\dots,e_{L-d}\}\) 与正整数
\(
s_1\mid s_2\mid\cdots\mid s_{L-d}
\)
使得
\[
M_{\ZZ}(P)=\bigoplus_{i=1}^{L-d}s_i\ZZ\,e_i,
\qquad
T_{\mathrm{prim}}\cong \bigoplus_{i=1}^{L-d}\ZZ/s_i\ZZ.
\]
模 \(p\) 约化后，
\[
\overline{K}_p=\bigoplus_{i=1}^{L-d}\FF_p\,\overline e_i,
\qquad
\overline{M}_p=\bigoplus_{i=1}^{L-d}\FF_p\,\overline{s_i}\,\overline e_i.
\]
因此
\[
\dim_{\FF_p}\frac{\overline{K}_p}{\overline{M}_p}
=
\#\{\,i:\ p\mid s_i\,\}
=
\dim_{\FF_p}\!\bigl(T_{\mathrm{prim}}\otimes\FF_p\bigr).
\]
结合 \(\overline{K}_p=\ker_{\FF_p}(\mathsf H_{N,L}\bmod p)\) 即得第 (2) 条。
最后，\(T_{\mathrm{prim}}\otimes\FF_p=0\) 当且仅当 \(p\nmid |T_{\mathrm{prim}}|\)，遂得失真素数集合的精确刻画。证毕。
\end{proof}

\endinput


\begin{proposition}[全部尺度上的 Hankel-null 稳定形态]\label{prop:pom-resonance-hankel-null-all-scales}
\leanverified{paper\_pom\_resonance\_hankel\_null\_all\_scales}
在定义 \ref{def:pom-resonance-hankel-null-multiples} 的设定下，对任意整数 \(m\ge d_q+1\)，令
\[
H^{(m)}_{i+1,j+1}:=a_{i+j}\qquad(0\le i\le m-1,\ 0\le j\le L-1),
\]
并记
\(
\mathrm{NullModes}^{(m)}(q):=\ker_{\ZZ}\bigl((H^{(m)})^{\mathsf T}\bigr)\subset\ZZ^{m}
\)。
定义
\[
\mathcal{M}^{(m)}(P_q):=
\Bigl\{\operatorname{coeff}_{<m}\!\bigl(Q(\lambda)P_q(\lambda)\bigr)\in\ZZ^m:\ 
Q\in\ZZ[\lambda],\ \deg Q<m-d_q\Bigr\}\subset\ZZ^m.
\]
若 \(L\ge m\)，则对一切 \(m\ge d_q+1\) 有
\[
\boxed{\;\mathrm{NullModes}^{(m)}(q)=\mathcal{M}^{(m)}(P_q),\qquad \mathrm{rank}\,\mathrm{NullModes}^{(m)}(q)=m-d_q.\;}
\]
\end{proposition}
\begin{proof}
对序列 \((a_t)\) 再次应用定理 \ref{thm:pom-hankel-syndrome-module-kernel-equals-multiples}，把其中的 \(n\) 替换为 \(m\) 即得方块 Hankel 的结论。
当 \(L\ge m\) 时，\((H^{(m)})^{\mathsf T}\) 的核条件与方块 Hankel 的核条件等价，从而得到所示恒等式。
秩公式由推论 \ref{cor:pom-hankel-syndrome-module-rank-and-generators}（取 \(n=m\)、\(d=d_q\)）推出。证毕。
\end{proof}

\begin{conjecture}[原始商消失与全尺度主理想性]\label{conj:pom-resonance-hankel-null-principal-all-q}
沿用定义 \ref{def:pom-resonance-hankel-null-multiples} 的记号。
对任意 \(q\ge 2\)，设 \(P_q\) 为相应的首一最小湮灭多项式，并定义原始商
\[
\Delta_{\mathrm{prim}}(q):=\dim\Bigl(\ker_{\ZZ}(H^{\mathsf T})\big/\mathcal{M}(P_q)\Bigr)
\]
（在可比的 Hankel 审计窗口与同型截断口径下）。
则对所有 \(q\ge 2\) 有 \(\Delta_{\mathrm{prim}}(q)=0\)，并且对所有 \(m\ge d_q+1\) 有
\[
\ker_{\ZZ}\bigl((H^{(m)})^{\mathsf T}\bigr)=\mathcal{M}^{(m)}(P_q)\subset\ZZ^m.
\]
该猜想在共振窗 \(q=9,\dots,17\) 已由表 \ref{tab:fold_collision_resonance_nullmodes_primitive_quotient_q9_17} 的精确倍乘分解与 \(\Delta_{\mathrm{prim}}(q)=0\) 的审计输出一致支持。
\end{conjecture}

\endinput
